% GENERATED FILE. Edit canonical lessons, then run npm run generate:books.
% canonical-source-hash: fnv1a64:e853157f40507e95
% canonical-lessons: KA-C166-ulu, KA-C166-cimukisu, KA-C166-tottikku, KA-C166-soru, KA-C166-telu, KA-R166-first-pass-words-from-chapters-158-161, KA-R166-second-pass-words-from-chapters-162-166

\chapter{To plough, To sprinkle, To drip, To leak, To float}
\label{ch:ka-to-plough-to-sprinkle-to-drip-to-leak-to-float}
\hlchaptermodality{166}

\begin{chapteropening}
\textbf{By the end of this chapter:} \emph{I can say to plough, to sprinkle, to drip, to leak and to float.}

Complete the last lesson of chapter 166: I can say to plough, to sprinkle, to drip, to leak and to float.
\end{chapteropening}

\section[\texorpdfstring{\kn{ಉಳು} (uḷu)}{ಉಳು (uḷu)}]{\kn{ಉಳು} (uḷu) — to plough}

\label{lesson:KA-C166-ulu}



\begin{quote}
Before the new one: say the Kannada for to remain, then the Kannada for to heat.
\end{quote}

\subsection*{You'll want to know: \kn{ಉಳು}}
\textbf{\kn{ಉಳು}} — \emph{uḷu} — \textquotedblleft{}to plough\textquotedblright{}.

\begin{grammarlens}[title={What you've built}]
One more word for this chapter.
\end{grammarlens}

\subsection*{Your turn}
\begin{itemize}
\raggedright
  \item \emph{Say it:} \emph{uḷu}
  \item \emph{Say it:} \emph{uḷu}, once more
  \item \emph{Say it:} say \emph{uḷu}
  \item \emph{Recall:} say the Kannada for to remain, then the Kannada for to heat, then say \emph{uḷu} again
\end{itemize}

\begin{tcolorbox}[breakable,colback=teal!4,colframe=teal!35!black,title={Before you move on}]
What does \kn{ಉಳು} mean? (\textquotedblleft{}To plough\textquotedblright{}.) Say it once more.
\end{tcolorbox}
\section[\texorpdfstring{\kn{ಚಿಮುಕಿಸು} (cimukisu)}{ಚಿಮುಕಿಸು (cimukisu)}]{\kn{ಚಿಮುಕಿಸು} (cimukisu) — to sprinkle}

\label{lesson:KA-C166-cimukisu}



\begin{quote}
Before the new one: say the Kannada for to heat, then the Kannada for to plough.
\end{quote}

\subsection*{You'll want to know: \kn{ಚಿಮುಕಿಸು}}
\textbf{\kn{ಚಿಮುಕಿಸು}} — \emph{cimukisu} — \textquotedblleft{}to sprinkle\textquotedblright{}.

\begin{grammarlens}[title={What you've built}]
One more word for this chapter.
\end{grammarlens}

\subsection*{Your turn}
\begin{itemize}
\raggedright
  \item \emph{Say it:} \emph{cimukisu}
  \item \emph{Say it:} \emph{cimukisu}, once more
  \item \emph{Say it:} say \emph{cimukisu}
  \item \emph{Recall:} say the Kannada for to heat, then the Kannada for to plough, then say \emph{cimukisu} again
\end{itemize}

\begin{tcolorbox}[breakable,colback=teal!4,colframe=teal!35!black,title={Before you move on}]
What does \kn{ಚಿಮುಕಿಸು} mean? (\textquotedblleft{}To sprinkle\textquotedblright{}.) Say it once more.
\end{tcolorbox}
\section[\texorpdfstring{\kn{ತೊಟ್ಟಿಕ್ಕು} (toṭṭikku)}{ತೊಟ್ಟಿಕ್ಕು (toṭṭikku)}]{\kn{ತೊಟ್ಟಿಕ್ಕು} (toṭṭikku) — to drip}

\label{lesson:KA-C166-tottikku}



\begin{quote}
Before the new one: say the Kannada for to plough, then the Kannada for to sprinkle.
\end{quote}

\subsection*{You'll want to know: \kn{ತೊಟ್ಟಿಕ್ಕು}}
\textbf{\kn{ತೊಟ್ಟಿಕ್ಕು}} — \emph{toṭṭikku} — \textquotedblleft{}to drip\textquotedblright{}.

\begin{grammarlens}[title={What you've built}]
One more word for this chapter.
\end{grammarlens}

\subsection*{Your turn}
\begin{itemize}
\raggedright
  \item \emph{Say it:} \emph{toṭṭikku}
  \item \emph{Say it:} \emph{toṭṭikku}, once more
  \item \emph{Say it:} say \emph{toṭṭikku}
  \item \emph{Recall:} say the Kannada for to plough, then the Kannada for to sprinkle, then say \emph{toṭṭikku} again
\end{itemize}

\begin{tcolorbox}[breakable,colback=teal!4,colframe=teal!35!black,title={Before you move on}]
What does \kn{ತೊಟ್ಟಿಕ್ಕು} mean? (\textquotedblleft{}To drip\textquotedblright{}.) Say it once more.
\end{tcolorbox}
\section[\texorpdfstring{\kn{ಸೋರು} (sōru)}{ಸೋರು (sōru)}]{\kn{ಸೋರು} (sōru) — to leak}

\label{lesson:KA-C166-soru}



\begin{quote}
Before the new one: say the Kannada for to sprinkle, then the Kannada for to drip.
\end{quote}

\subsection*{You'll want to know: \kn{ಸೋರು}}
\textbf{\kn{ಸೋರು}} — \emph{sōru} — \textquotedblleft{}to leak\textquotedblright{}.

\begin{grammarlens}[title={What you've built}]
One more word for this chapter.
\end{grammarlens}

\subsection*{Your turn}
\begin{itemize}
\raggedright
  \item \emph{Say it:} \emph{sōru}
  \item \emph{Say it:} \emph{sōru}, once more
  \item \emph{Say it:} say \emph{sōru}
  \item \emph{Recall:} say the Kannada for to sprinkle, then the Kannada for to drip, then say \emph{sōru} again
\end{itemize}

\begin{tcolorbox}[breakable,colback=teal!4,colframe=teal!35!black,title={Before you move on}]
What does \kn{ಸೋರು} mean? (\textquotedblleft{}To leak\textquotedblright{}.) Say it once more.
\end{tcolorbox}
\section[\texorpdfstring{\kn{ತೇಲು} (tēlu)}{ತೇಲು (tēlu)}]{\kn{ತೇಲು} (tēlu) — to float}

\label{lesson:KA-C166-telu}



\begin{quote}
Before the new one: say the Kannada for to drip, then the Kannada for to leak.
\end{quote}

\subsection*{You'll want to know: \kn{ತೇಲು}}
\textbf{\kn{ತೇಲು}} — \emph{tēlu} — \textquotedblleft{}to float\textquotedblright{}.

\begin{grammarlens}[title={What you've built}]
One more word for this chapter.
\end{grammarlens}

\subsection*{Your turn}
\begin{itemize}
\raggedright
  \item \emph{Say it:} \emph{tēlu}
  \item \emph{Say it:} \emph{tēlu}, once more
  \item \emph{Say it:} say \emph{tēlu}
  \item \emph{Recall:} say the Kannada for to drip, then the Kannada for to leak, then say \emph{tēlu} again
\end{itemize}

\begin{tcolorbox}[breakable,colback=teal!4,colframe=teal!35!black,title={Before you move on}]
What does \kn{ತೇಲು} mean? (\textquotedblleft{}To float\textquotedblright{}.) Say it once more.
\end{tcolorbox}
\section[(dialogue)]{Words from chapters 158-161 — a first pass}

\label{lesson:KA-R166-first-pass-words-from-chapters-158-161}



\begin{quote}
Say the words for to leak and to float.
\end{quote}

\subsection*{Your turn}
Say these aloud:

\begin{itemize}
\raggedright
  \item to boil, to cook, to buy, to pay, to bathe
  \item to peel, to grind, to serve, to taste, to chew
  \item to bend down, to carry, to exercise, to sneeze, to shiver
  \item to stick, to inform, to meet, to reserve, to cancel
\end{itemize}

\begin{tcolorbox}[breakable,colback=teal!4,colframe=teal!35!black,title={Before you move on}]
To boil? (\textbf{\kn{ಬೇಯಿಸು}}, \emph{bēyisu}.) To cook? (\textbf{\kn{ಅಡುಗೆ} \kn{ಮಾಡು}}, \emph{aḍuge māḍu}.) To buy? (\textbf{\kn{ಖರೀದಿಸು}}, \emph{kharīdisu}.) To pay? (\textbf{\kn{ಪಾವತಿಸು}}, \emph{pāvatisu}.) To bathe? (\textbf{\kn{ಸ್ನಾನ} \kn{ಮಾಡು}}, \emph{snāna māḍu}.) To peel? (\textbf{\kn{ಸುಲಿ}}, \emph{suli}.) To grind? (\textbf{\kn{ಅರೆ}}, \emph{are}.) To serve? (\textbf{\kn{ಬಡಿಸು}}, \emph{baḍisu}.) To taste? (\textbf{\kn{ರುಚಿ} \kn{ನೋಡು}}, \emph{ruci nōḍu}.) To chew? (\textbf{\kn{ಅಗಿ}}, \emph{agi}.) To bend down? (\textbf{\kn{ಬಗ್ಗು}}, \emph{baggu}.) To carry? (\textbf{\kn{ಹೊರು}}, \emph{horu}.) To exercise? (\textbf{\kn{ವ್ಯಾಯಾಮ} \kn{ಮಾಡು}}, \emph{vyāyāma māḍu}.) To sneeze? (\textbf{\kn{ಸೀನು}}, \emph{sīnu}.) To shiver? (\textbf{\kn{ನಡುಗು}}, \emph{naḍugu}.) To stick? (\textbf{\kn{ಅಂಟಿಸು}}, \emph{aṇṭisu}.) To inform? (\textbf{\kn{ತಿಳಿಸು}}, \emph{tiḷisu}.) To meet? (\textbf{\kn{ಭೇಟಿ} \kn{ಮಾಡು}}, \emph{bhēṭi māḍu}.) To reserve? (\textbf{\kn{ಕಾಯ್ದಿರಿಸು}}, \emph{kāydirisu}.) To cancel? (\textbf{\kn{ರದ್ದು} \kn{ಮಾಡು}}, \emph{raddu māḍu}.)
\end{tcolorbox}
\section[(dialogue)]{Words from chapters 162-166 — a second pass}

\label{lesson:KA-R166-second-pass-words-from-chapters-162-166}



\begin{quote}
Say the words for to correct and to copy.
\end{quote}

\subsection*{Your turn}
Say these aloud:

\begin{itemize}
\raggedright
  \item to correct, to copy, to warn, to take part, to transfer
  \item to doubt, to put up with, to cheat, to get along, to plant
  \item to upload, to click, to delete, to switch on, to switch off
  \item to steal, to recall, to roll, to remain, to heat
  \item to plough, to sprinkle, to drip, to leak, to float
\end{itemize}

\begin{tcolorbox}[breakable,colback=teal!4,colframe=teal!35!black,title={Before you move on}]
To correct? (\textbf{\kn{ತಿದ್ದು}}, \emph{tiddu}.) To copy? (\textbf{\kn{ನಕಲು} \kn{ಮಾಡು}}, \emph{nakalu māḍu}.) To warn? (\textbf{\kn{ಎಚ್ಚರಿಸು}}, \emph{eccarisu}.) To take part? (\textbf{\kn{ಭಾಗವಹಿಸು}}, \emph{bhāgavahisu}.) To transfer? (\textbf{\kn{ವರ್ಗಾಯಿಸು}}, \emph{vargāyisu}.) To doubt? (\textbf{\kn{ಅನುಮಾನಪಡು}}, \emph{anumānapaḍu}.) To put up with? (\textbf{\kn{ಸಹಿಸು}}, \emph{sahisu}.) To cheat? (\textbf{\kn{ಮೋಸ} \kn{ಮಾಡು}}, \emph{mōsa māḍu}.) To get along? (\textbf{\kn{ಹೊಂದಿಕೊ}}, \emph{hondiko}.) To plant? (\textbf{\kn{ನೆಡು}}, \emph{neḍu}.) To upload? (\textbf{\kn{ಅಪ್ಲೋಡ್} \kn{ಮಾಡು}}, \emph{aplōḍ māḍu}.) To click? (\textbf{\kn{ಕ್ಲಿಕ್} \kn{ಮಾಡು}}, \emph{klik māḍu}.) To delete? (\textbf{\kn{ಅಳಿಸು}}, \emph{aḷisu}.) To switch on? (\textbf{\kn{ಆನ್} \kn{ಮಾಡು}}, \emph{ān māḍu}.) To switch off? (\textbf{\kn{ಆಫ್} \kn{ಮಾಡು}}, \emph{āph māḍu}.) To steal? (\textbf{\kn{ಕದಿ}}, \emph{kadi}.) To recall? (\textbf{\kn{ನೆನಪಿಸಿಕೊಳ್ಳು}}, \emph{nenapisikoḷḷu}.) To roll? (\textbf{\kn{ಉರುಳು}}, \emph{uruḷu}.) To remain? (\textbf{\kn{ಉಳಿ}}, \emph{uḷi}.) To heat? (\textbf{\kn{ಕಾಯಿಸು}}, \emph{kāyisu}.) To plough? (\textbf{\kn{ಉಳು}}, \emph{uḷu}.) To sprinkle? (\textbf{\kn{ಚಿಮುಕಿಸು}}, \emph{cimukisu}.) To drip? (\textbf{\kn{ತೊಟ್ಟಿಕ್ಕು}}, \emph{toṭṭikku}.) To leak? (\textbf{\kn{ಸೋರು}}, \emph{sōru}.) To float? (\textbf{\kn{ತೇಲು}}, \emph{tēlu}.)
\end{tcolorbox}
